% Options for packages loaded elsewhere
\PassOptionsToPackage{unicode}{hyperref}
\PassOptionsToPackage{hyphens}{url}
\PassOptionsToPackage{dvipsnames,svgnames,x11names}{xcolor}
\documentclass[
  11pt,
]{article}
\usepackage{xcolor}
\usepackage[margin=25mm]{geometry}
\usepackage{amsmath,amssymb}
\setcounter{secnumdepth}{-\maxdimen} % remove section numbering
\usepackage{iftex}
\ifPDFTeX
  \usepackage[T1]{fontenc}
  \usepackage[utf8]{inputenc}
  \usepackage{textcomp} % provide euro and other symbols
\else % if luatex or xetex
  \usepackage{unicode-math} % this also loads fontspec
  \defaultfontfeatures{Scale=MatchLowercase}
  \defaultfontfeatures[\rmfamily]{Ligatures=TeX,Scale=1}
\fi
\usepackage{lmodern}
\ifPDFTeX\else
  % xetex/luatex font selection
\fi
% Use upquote if available, for straight quotes in verbatim environments
\IfFileExists{upquote.sty}{\usepackage{upquote}}{}
\IfFileExists{microtype.sty}{% use microtype if available
  \usepackage[]{microtype}
  \UseMicrotypeSet[protrusion]{basicmath} % disable protrusion for tt fonts
}{}
\makeatletter
\@ifundefined{KOMAClassName}{% if non-KOMA class
  \IfFileExists{parskip.sty}{%
    \usepackage{parskip}
  }{% else
    \setlength{\parindent}{0pt}
    \setlength{\parskip}{6pt plus 2pt minus 1pt}}
}{% if KOMA class
  \KOMAoptions{parskip=half}}
\makeatother
\usepackage{listings}
\newcommand{\passthrough}[1]{#1}
\lstset{defaultdialect=[5.3]Lua}
\lstset{defaultdialect=[x86masm]Assembler}
\usepackage{longtable,booktabs,array}
\usepackage{calc} % for calculating minipage widths
% Correct order of tables after \paragraph or \subparagraph
\usepackage{etoolbox}
\makeatletter
\patchcmd\longtable{\par}{\if@noskipsec\mbox{}\fi\par}{}{}
\makeatother
% Allow footnotes in longtable head/foot
\IfFileExists{footnotehyper.sty}{\usepackage{footnotehyper}}{\usepackage{footnote}}
\makesavenoteenv{longtable}
\setlength{\emergencystretch}{3em} % prevent overfull lines
\providecommand{\tightlist}{%
  \setlength{\itemsep}{0pt}\setlength{\parskip}{0pt}}
\emergencystretch=6em
\usepackage{graphicx}
\usepackage{xurl}
\usepackage{pdflscape}
\newsavebox{\PfpMathBox}
\newcommand{\PfpDisplayMath}[1]{%
  \sbox{\PfpMathBox}{$\displaystyle #1$}%
  \ifdim\wd\PfpMathBox>\linewidth
    \resizebox{\linewidth}{!}{\usebox{\PfpMathBox}}%
  \else\usebox{\PfpMathBox}\fi}
\AtBeginDocument{\Urlmuskip=0mu plus 2mu}
\lstset{breaklines=true,basicstyle=\small\ttfamily,columns=fullflexible,keepspaces=true}
\setlength{\tabcolsep}{4pt}
\usepackage{iftex}
\ifXeTeX
  \usepackage{newunicodechar}
  \makeatletter
  \newcommand{\DeclareUnicodeCharacter}[2]{%
    \begingroup
    \uccode`\~="#1\relax
    \uppercase{\def\temp{~}}%
    \expandafter\endgroup
    \expandafter\newunicodechar\expandafter{\temp}{#2}%
  }
  \makeatother
\fi
\ifdefined\DeclareUnicodeCharacter
\DeclareUnicodeCharacter{2192}{\ensuremath{\to}}
\DeclareUnicodeCharacter{2190}{\ensuremath{\leftarrow}}
\DeclareUnicodeCharacter{2194}{\ensuremath{\leftrightarrow}}
\DeclareUnicodeCharacter{21D2}{\ensuremath{\Rightarrow}}
\DeclareUnicodeCharacter{21A6}{\ensuremath{\mapsto}}
\DeclareUnicodeCharacter{2212}{\ensuremath{-}}
\DeclareUnicodeCharacter{2264}{\ensuremath{\leq}}
\DeclareUnicodeCharacter{2265}{\ensuremath{\geq}}
\DeclareUnicodeCharacter{2260}{\ensuremath{\neq}}
\DeclareUnicodeCharacter{2113}{\ensuremath{\ell}}
\DeclareUnicodeCharacter{03A3}{\ensuremath{\Sigma}}
\DeclareUnicodeCharacter{03A6}{\ensuremath{\Phi}}
\DeclareUnicodeCharacter{0393}{\ensuremath{\Gamma}}
\DeclareUnicodeCharacter{221E}{\ensuremath{\infty}}
\DeclareUnicodeCharacter{2208}{\ensuremath{\in}}
\DeclareUnicodeCharacter{2209}{\ensuremath{\notin}}
\DeclareUnicodeCharacter{211A}{\ensuremath{\mathbb{Q}}}
\DeclareUnicodeCharacter{211D}{\ensuremath{\mathbb{R}}}
\DeclareUnicodeCharacter{2124}{\ensuremath{\mathbb{Z}}}
\DeclareUnicodeCharacter{2115}{\ensuremath{\mathbb{N}}}
\DeclareUnicodeCharacter{2200}{\ensuremath{\forall}}
\DeclareUnicodeCharacter{2203}{\ensuremath{\exists}}
\DeclareUnicodeCharacter{2286}{\ensuremath{\subseteq}}
\DeclareUnicodeCharacter{00D7}{\ensuremath{\times}}
\DeclareUnicodeCharacter{03B4}{\ensuremath{\delta}}
\DeclareUnicodeCharacter{2202}{\ensuremath{\partial}}
\DeclareUnicodeCharacter{03B2}{\ensuremath{\beta}}
\DeclareUnicodeCharacter{03B1}{\ensuremath{\alpha}}
\DeclareUnicodeCharacter{03C6}{\ensuremath{\phi}}
\DeclareUnicodeCharacter{03B3}{\ensuremath{\gamma}}
\DeclareUnicodeCharacter{00B9}{\ensuremath{^1}}
\DeclareUnicodeCharacter{00B2}{\ensuremath{^2}}
\DeclareUnicodeCharacter{2070}{\ensuremath{^0}}
\DeclareUnicodeCharacter{2080}{\ensuremath{_0}}
\DeclareUnicodeCharacter{2081}{\ensuremath{_1}}
\DeclareUnicodeCharacter{2082}{\ensuremath{_2}}
\fi
\usepackage{bookmark}
\IfFileExists{xurl.sty}{\usepackage{xurl}}{} % add URL line breaks if available
\urlstyle{same}
\hypersetup{
  pdftitle={PFP V: A Practitioner\textquotesingle s Guide to Checked Proposals},
  pdfauthor={Jeremy H. Carroll},
  colorlinks=true,
  linkcolor={Maroon},
  filecolor={Maroon},
  citecolor={Blue},
  urlcolor={Blue},
  pdfcreator={LaTeX via pandoc}}

\title{PFP V: A Practitioner\textquotesingle s Guide to Checked
Proposals}
\author{Jeremy H. Carroll}
\date{2026-10-23 · Version 0.2.0 · Author-review extension}

\begin{document}
\makeatletter
\begin{center}
{\LARGE \@title\par}
\vspace{0.8em}
{\normalsize \@author\par}
\vspace{0.5em}
{\small \@date\par}
\end{center}
\makeatother

\textbf{Code baseline:}
\nolinkurl{59dde95a7860063c25d75db37f40af556e9f28ee}

\subsection{Abstract}\label{abstract}

This instructional preprint presents a reproducible proposal-check
workflow on a finite directed graph with exact rational edge claims. A
three-edge inconsistency has a hand-checkable weighted minimum of three.
The worked session distinguishes minimizing the residual of unchanged
claims from admitting a consistent replacement, and follows both
refusals and restoration of a saved state. The companion provides an
exact Python enumeration route and a separately configured native route,
with explicit input, verification, and evidence boundaries. The
Landau-style examination and study guide extend the session through
counterexamples and an executable capstone. The contribution is a scoped
practitioner workflow; no new theorem or measured model-performance
result is claimed.

\subsection{Purpose and reader
outcome}\label{purpose-and-reader-outcome}

This guide takes one small inconsistent graph through a complete
session: declare the instance, propose a change, interpret a refusal,
admit a repair, check the result, and restore an earlier state. It is a
practitioner companion to the Proposal Fidelity Protocol (PFP) research
series and its legacy Symbol Game \nolinkurl{implementations.} It adds
no new theorem or measured model-performance result.

A \textbf{proposal} is a candidate supplied by a person, script, or
model. A \textbf{host} is the program that applies the declared checks.
A \textbf{receipt} records what a particular execution checked and
returned. These roles matter because a plausible proposal can be wrong,
while a successful check can establish less than its surrounding prose
suggests.

The first session uses a scripted proposer. No model account, API key,
or live conversation is needed. You should finish able to explain both
an accepted move and the limits of its acceptance. The companion
\href{_lev_landau_theoretical_minimum_quiz/quiz.md}{Landau quiz} has 32
questions in eight levels; its
\href{_lev_landau_theoretical_minimum_quiz/study_guide.md}{study guide
and instructor key} supplies worked answers, refusal cases, and a
capstone replay.

\subsection{1. Declare the instance before solving
it}\label{declare-the-instance-before-solving-it}

Suppose three events have proposed offsets on one linear time axis:

\begin{itemize}
\tightlist
\item
  Event b is one unit after event a.
\item
  Event c is one unit after event b.
\item
  Event a is one unit after event c.
\end{itemize}

This is a mathematical offset example. There is no wraparound clock or
recurring-day convention. Declare those conventions differently and you
have a different problem.

Represent the events by vertices a, b, c and the statements by directed
edges. The ordered claim vector is

\begin{center}\PfpDisplayMath{
c=(c_{ab},c_{bc},c_{ca})=(1,1,1).
}\end{center}

A \textbf{vertex potential} assigns a number
\(f_a,\allowbreak f_b,\allowbreak f_c\) to the vertices. Its edge
differences are

\begin{center}\PfpDisplayMath{
df=(f_b-f_a,\ f_c-f_b,\ f_a-f_c).
}\end{center}

The differences always sum to zero: each vertex value appears once with
each sign. The supplied claims sum to three. Consequently, no potential
satisfies all three unchanged claims.

Declare rational coefficients for claims and potentials, ordinary
additive differences, and positive exact edge weights
\(w=(1,\allowbreak 2,\allowbreak 3)\). A weight is the penalty for
leaving one unit of discrepancy on that edge. It is a modeling choice,
not evidence that one event is more truthful. The complete instance
therefore includes the vertices, directed edges, edge order, exact
claims, weights, units, and permitted changes. The graph has no declared
two-dimensional cells. Its directed arrows alone do not attach a filled
triangle.

\subsection{2. Separate a residual from a changed
claim}\label{separate-a-residual-from-a-changed-claim}

For a chosen potential, the \textbf{residual} is \(r=c-df\). Its
weighted discrepancy is

\begin{center}\PfpDisplayMath{
D_c(f)=|r_{ab}|+2|r_{bc}|+3|r_{ca}|.
}\end{center}

The \textbf{quotient minimum} is the smallest such discrepancy over all
potentials:

\begin{center}\PfpDisplayMath{
\Phi_w(c)=\min_f D_c(f).
}\end{center}

The notation distinguishes the score at one chosen potential from the
best possible score for unchanged claims. A poor potential can make
\(D_c(f)\) larger without changing \(\Phi_w(c)\).

Take \(f=(0,\allowbreak -2,\allowbreak -1)\). Then
\(df=(-2,\allowbreak 1,\allowbreak 1)\), so
\(r=(3,\allowbreak 0,\allowbreak 0)\) and \(D_c(f)=3\). For every
potential, the residual still sums to three. Because all weights are at
least one,

\begin{center}\PfpDisplayMath{
D_c(f)\geq |r_{ab}|+|r_{bc}|+|r_{ca}|
\geq |r_{ab}+r_{bc}+r_{ca}|=3.
}\end{center}

The displayed potential attains this lower bound. This short argument
establishes the minimum for this instance without requiring trust in a
solver.

Two operations now have different meanings. Choosing that potential
concentrates the original discrepancy on one edge; the original cycle
sum remains three. Replacing the claims themselves by
\(c'=(-2,\allowbreak 1,\allowbreak 1)\) produces a new consistent
instance: \(c'=df\) and \(\Phi_w(c')=0\). The governor demonstration
performs this second operation through
\nolinkurl{step_cut(proposed_claims=...)}. Acceptance says the changed
claims satisfy the declared graph check. It does not establish that the
changed event statement is factually correct.

\begin{longtable}[]{@{}
  >{\raggedright\arraybackslash}p{(\linewidth - 6\tabcolsep) * \real{0.2499}}
  >{\raggedright\arraybackslash}p{(\linewidth - 6\tabcolsep) * \real{0.2499}}
  >{\raggedright\arraybackslash}p{(\linewidth - 6\tabcolsep) * \real{0.2499}}
  >{\raggedright\arraybackslash}p{(\linewidth - 6\tabcolsep) * \real{0.2499}}@{}}
\toprule\noalign{}
\begin{minipage}[b]{\linewidth}\raggedright
Claims
\end{minipage} & \begin{minipage}[b]{\linewidth}\raggedright
Chosen potential
\end{minipage} & \begin{minipage}[b]{\linewidth}\raggedright
Score at that potential
\end{minipage} & \begin{minipage}[b]{\linewidth}\raggedright
Minimum for those claims
\end{minipage} \\
\midrule\noalign{}
\endhead
\bottomrule\noalign{}
\endlastfoot
\((1,\allowbreak 1,\allowbreak 1)\) &
\((0,\allowbreak 0,\allowbreak 0)\) & 6 & 3 \\
\((1,\allowbreak 1,\allowbreak 1)\) &
\((0,\allowbreak -2,\allowbreak -1)\) & 3 & 3 \\
\((-2,\allowbreak 1,\allowbreak 1)\) &
\((0,\allowbreak 0,\allowbreak 0)\) & 7 & 0 \\
\((-2,\allowbreak 1,\allowbreak 1)\) &
\((0,\allowbreak -2,\allowbreak -1)\) & 0 & 0 \\
\end{longtable}

Adding the same rational constant to all three potential values changes
no edge difference. The displayed potential fixes \(f_a=0\) for
convenience. The unique cheapest edge fixes the minimizing residual in
this instance; it does not make a numerical gauge offset observable.
With tied cheapest weights, different minimizing residuals can tie as
well.

\subsection{3. Choose the executable
path}\label{choose-the-executable-path}

The supported tutorial in this package is
\href{code/quickstart.py}{\nolinkurl{code/quickstart.py}}. Its Python
wrapper sends exact graph inputs to the pinned Rust calculator. The
repository also contains a separate historical root
\nolinkurl{quickstart.py}; the commands here intentionally name the
package file.

Use Python 3.10 or newer. The Python companion is self-contained: it
carries copies of the native graph adapter
\href{code/native_functional.py}{\nolinkurl{code/native_functional.py}}
and the Paper 004 graph governor
\href{code/l1_symbol_game_governor.py}{\nolinkurl{code/l1_symbol_game_governor.py}}.
Each header records the source path, commit, source-byte SHA-256, and
local changes. The adapter differs from its baseline only in its header
and import order. The governor has documented patches to
\passthrough{\lstinline!compute\_phi1!},
\passthrough{\lstinline!step\_assert!},
\passthrough{\lstinline!step\_cut!}, and
\passthrough{\lstinline!step\_verify!}: input normalization and carrier
checks, exact-value admission, complete replacement claims, atomic
rollback on a projection error, scoped edge verification, and
native-refusal propagation.
\nolinkurl{code/tests/test_vendored_sources.py} compares the remaining
executable structure with the recorded source objects; its explicitly
excluded patched functions and methods have behavioral regression checks
in \nolinkurl{code/tests/test_manual_contracts.py} and the earlier
boundary tests. Source-object comparisons skip with a reason when the
Git objects are absent.

The native calculator itself is not part of this package. For the
baseline route, use a separate checkout or worktree at the recorded code
baseline, confirm its revision with
\passthrough{\lstinline!git rev-parse HEAD!}, then build
\nolinkurl{solve_graph_phi1} with that checkout's Rust workspace:

\begin{lstlisting}
cargo build --locked -j 1 --manifest-path omega_engine/crates_rs/Cargo.toml -p omega_functional --bin solve_graph_phi1
\end{lstlisting}

The following Windows PowerShell setup, run from the repository root,
assumes the reviewed executable already exists at the shown path:

\begin{lstlisting}
$env:OMEGA_FUNCTIONAL_GRAPH_BIN = (Resolve-Path omega_engine/crates_rs/target/debug/solve_graph_phi1.exe).Path
$env:OMEGA_FUNCTIONAL_GRAPH_SHA256 = (Get-FileHash -Algorithm SHA256 $env:OMEGA_FUNCTIONAL_GRAPH_BIN).Hash.ToLowerInvariant()
\end{lstlisting}

Use the actual output path if Cargo has a custom target directory. Other
platforms use their native executable name. The digest pins the
executable bytes used by the adapter; it is not an independent
attestation that those bytes were reproducibly built from the intended
source.

Now, from this package's directory, run (on Windows,
\passthrough{\lstinline!py -3.12!} in place of
\passthrough{\lstinline!python!}):

\begin{lstlisting}
python code/quickstart.py --non-interactive
python code/play_demo.py
\end{lstlisting}

The repository's
\href{https://github.com/dbccpage/the_symbol_game/blob/52cb5457f6c485ba548cfd292f6012a941ed22cf/omega_engine_education/FUNCTIONAL_BOUNDARY.md}{functional-boundary
notes} describe the native setup at the baseline. The transport
implementation is
\href{code/native_functional.py}{\nolinkurl{native_functional.py}}. A
missing executable, missing or mismatched digest, native refusal, or
malformed native response prevents an accepted native result. Treat that
outcome as an environment or evidence issue to resolve before relying on
the host calculation.

\subsubsection{Offline arithmetic when the native calculator is
unavailable}\label{offline-arithmetic-when-the-native-calculator-is-unavailable}

You can still understand and check the finite argument. Run this
ordinary Python snippet, or evaluate it by hand:

\begin{lstlisting}[language=Python]
from fractions import Fraction as Q

claims = (Q(1), Q(1), Q(1))
potential = (Q(0), Q(-2), Q(-1))
edge_differences = (
    potential[1] - potential[0],
    potential[2] - potential[1],
    potential[0] - potential[2],
)
residual = tuple(c - d for c, d in zip(claims, edge_differences))
assert edge_differences == (Q(-2), Q(1), Q(1))
assert residual == (Q(3), Q(0), Q(0))
assert sum(edge_differences) == 0
assert sum(residual) == 3
\end{lstlisting}

This is an arithmetic worksheet. It neither calls the native adapter nor
emits a canonical functional receipt. Its lower-bound explanation is the
argument in Section 2. The package test
\nolinkurl{code/tests/test_offline_worksheet.py} runs this block as
printed. Record a native run as unavailable until the required
calculator can execute.

\subsection{4. Follow the complete
session}\label{follow-the-complete-session}

\href{code/play_demo.py}{play\_demo.py} imports the package's copy of
the \href{code/l1_symbol_game_governor.py}{graph governor}. It uses the
manuscript's weights (1,2,3) throughout. The all-zero potential costs
six; the original quotient minimum is three. The session prints and
returns its selected \nolinkurl{calculation_path}
(\passthrough{\lstinline!native!} or \nolinkurl{python-enumeration}) and
its weights. This path label describes the calculator selection; it is
not a source-to-binary attestation or a native receipt.

When \nolinkurl{OMEGA_FUNCTIONAL_GRAPH_BIN} is set, the governor asks
that native calculator, and any refusal or error it returns stops the
step; no Python value is substituted. When it is not set, the governor
computes the quotient minimum by exact Python enumeration of
spanning-tree potentials. Without a configured calculator the session
therefore still runs, but its printed values are Python enumeration
results, not native results. Edge weights must be exact integers or
fractions on both paths; floating-point and Boolean weights are refused.
The package test \nolinkurl{code/tests/test_play_demo.py} checks the
table below on the enumeration path, and again through the native
calculator when pins are configured. Enumeration can grow
combinatorially with the graph; this small-session route is not a
general scalability claim or an implementation of a proposal-attempt
budget.

\begin{longtable}[]{@{}
  >{\raggedright\arraybackslash}p{(\linewidth - 4\tabcolsep) * \real{0.3332}}
  >{\raggedright\arraybackslash}p{(\linewidth - 4\tabcolsep) * \real{0.3332}}
  >{\raggedright\arraybackslash}p{(\linewidth - 4\tabcolsep) * \real{0.3332}}@{}}
\toprule\noalign{}
\begin{minipage}[b]{\linewidth}\raggedright
Step
\end{minipage} & \begin{minipage}[b]{\linewidth}\raggedright
Proposed action
\end{minipage} & \begin{minipage}[b]{\linewidth}\raggedright
Expected observation
\end{minipage} \\
\midrule\noalign{}
\endhead
\bottomrule\noalign{}
\endlastfoot
Declare & Create vertices a, b, c and weighted edges & The host has a
finite declared carrier \\
ASSERT & Install claims \(1,\allowbreak 1,\allowbreak 1\) & The original
quotient minimum is 3 \\
CUT & Propose replacement \(5,\allowbreak 0,\allowbreak 0\) &
\nolinkurl{DefectIncrease}, not admitted \\
CUT & Propose replacement \(-2,\allowbreak 1,\allowbreak 1\) &
\passthrough{\lstinline!ExactRepair!}, admitted, 3 to 0 \\
VERIFY & Request \nolinkurl{ZERO_OBSTRUCTION} & The host checks its
stored claims and admits the check \\
RETRACT & Restore the saved post-ASSERT snapshot & The original claims
and minimum 3 return \\
\end{longtable}

The bad proposal sums to five. Its optimum is five by the same
lower-bound argument, with all discrepancy placeable on the weight-one
edge. The host refuses an increase from three to five. A refused
proposal does not replace the active claims.

The good proposal changes the cheapest edge from one to minus two. It
supplies a consistent replacement. The demonstration saves a snapshot
before either CUT and uses that snapshot's actual identifier when
retracting. Do not guess a snapshot identifier from a printed step
count.

ASSERT installs the starting problem; it is not a claim that every
initialization step decreases discrepancy. The governor's CUT checks
permit \nolinkurl{non-increase.} A tied score can therefore return
\nolinkurl{IdempotentStationary}. Strict progress and finite termination
of an arbitrary search require additional conditions; repeating an
admitted stationary move is not a solution strategy.

\nolinkurl{IdempotentStationary} describes an unchanged checked score,
not necessarily unchanged claims. For example, replacement claims
(2,1,0) still have minimum three and are admitted with that status.
\passthrough{\lstinline!RETRACT!} can deliberately raise the score when
restoring the original problem, so the CUT rule is not a monotonicity
theorem for the whole session.

\nolinkurl{ZERO_OBSTRUCTION} asks whether some potential fits the stored
claims. \passthrough{\lstinline!ZERO\_DEFECT!} asks whether the
currently stored potential fits them. Consistent claims (-2,1,1) with
the all-zero potential give a minimum of zero but a current discrepancy
of seven. Similarly, a successful
\passthrough{\lstinline!EDGE\_MATCH:b,c!} checks only that declared
edge; it does not discharge the original global cycle obstruction.
Undeclared or malformed \passthrough{\lstinline!EDGE\_MATCH!}
obligations now return \nolinkurl{VerificationFailed} rather than
verifying against default zeros.

\subsection{5. Read refusals as
information}\label{read-refusals-as-information}

A refusal identifies a failed check. First determine which check ran and
which input caused it. Then change the relevant proposal or declare a
new instance when the intended problem requires different rules.

\begin{longtable}[]{@{}
  >{\raggedright\arraybackslash}p{(\linewidth - 4\tabcolsep) * \real{0.3332}}
  >{\raggedright\arraybackslash}p{(\linewidth - 4\tabcolsep) * \real{0.3332}}
  >{\raggedright\arraybackslash}p{(\linewidth - 4\tabcolsep) * \real{0.3332}}@{}}
\toprule\noalign{}
\begin{minipage}[b]{\linewidth}\raggedright
Observed result
\end{minipage} & \begin{minipage}[b]{\linewidth}\raggedright
Meaning in this implementation
\end{minipage} & \begin{minipage}[b]{\linewidth}\raggedright
Next operator action
\end{minipage} \\
\midrule\noalign{}
\endhead
\bottomrule\noalign{}
\endlastfoot
\nolinkurl{CarrierViolation} & A relational CUT endpoint or section CUT
key is outside the declared vertices & Correct the endpoint or
explicitly formulate another carrier \\
\nolinkurl{MoveViolation} & A proposed edge is undeclared, or one CUT
supplies both section and replacement-claim payloads & Use one permitted
operation on the declared graph \\
\nolinkurl{IncompleteProposal} & A replacement CUT omits an edge, or
ASSERT/CUT supplies a non-exact mathematical value & Give exact integers
or fractions; a replacement CUT must explicitly cover every edge \\
\nolinkurl{InvariantViolation} & A prohibited input, such as a
self-edge, reached the gate & Inspect the exact failed input
condition \\
\nolinkurl{DefectIncrease} & The proposed score exceeds the current
score under the selected CUT mode & Compute why it increased before
proposing another change \\
\nolinkurl{IdempotentStationary} & A permitted move left the checked
score unchanged & Identify a new argument or record the lack of
progress \\
\nolinkurl{VerificationFailed} & The requested obligation did not pass &
Inspect that obligation, rather than treating a different successful
check as a substitute \\
\end{longtable}

The native transport can also fail before a governor verdict exists.
Keep the transport exception and native response, if any, with the run
record. A missing calculation is not a zero discrepancy.

Malformed inputs to \passthrough{\lstinline!compute\_phi1!} can raise
\passthrough{\lstinline!ValueError!} or
\passthrough{\lstinline!TypeError!} before a verdict. Unknown claim
edges and duplicate carrier entries are refused on both routes;
supported \passthrough{\lstinline!Edge!} object keys are normalized
before either route scores them. At initialization, unasserted declared
edges have zero claims, an explicit legacy convention. A replacement CUT
still must name every edge; it cannot rely on that initialization
default. Section CUTs retain the legacy default of zero for omitted
declared vertices, but refuse undeclared vertices and non-exact supplied
values. The guide uses complete potentials to make that convention
visible. An ASSERT with an undeclared endpoint returns the legacy
\nolinkurl{InvariantViolation} status rather than
\nolinkurl{CarrierViolation}.

A refused replacement or a projection exception preserves the active
claims, section, and admitted snapshot history. The attempt counter may
advance, so unchanged active state does not mean an identical entire
runtime record. Preserve refused attempts separately from the admitted
history, which does not contain every refusal or VERIFY result.

Use integers and \passthrough{\lstinline!Fraction!} inputs for
mathematical values. The tutorial rejects floats and Boolean values at
its mathematical input boundary. In particular, Python's treatment of
\passthrough{\lstinline!True!} as an integer-like value must not turn it
into a declared numerical witness.

\subsection{6. Reproduce a result and inspect its
custody}\label{reproduce-a-result-and-inspect-its-custody}

A useful reproduction record names the source revision, exact input,
command, Python version, native executable digest, exit status, and
output. If files change during a run, identify which version each
observation belongs to. Preserve refused cases alongside admitted ones:
they show whether the gate \nolinkurl{discriminates.}

Run the package checks from this package's directory, one test file per
command:

\begin{lstlisting}
python -m pytest code/tests/test_offline_worksheet.py -q
python -m pytest code/tests/test_play_demo.py -q
python -m pytest code/tests/test_quickstart_native.py -q
python -m pytest code/tests/test_manual_contracts.py -q
\end{lstlisting}

The \href{code/tests/conftest.py}{fixture} consumes only two explicitly
supplied pins:

\begin{itemize}
\tightlist
\item
  Executable path: \nolinkurl{OMEGA_FUNCTIONAL_GRAPH_BIN}.
\item
  Executable digest: \nolinkurl{OMEGA_FUNCTIONAL_GRAPH_SHA256}.
\end{itemize}

It discovers no checkout, builds no program, and creates no pin. Missing
configuration skips, with that reason, every test that requests native
execution. A malformed digest, unreadable configured executable, or
digest mismatch fails those tests with an explicit message. The
worksheet, the offline session, the refusal, exact-input, parser, and
scaffold checks, and the vendored-source header checks still run without
native \nolinkurl{configuration.} The commands are replay instructions,
not a reported fresh pass.

The existing \href{code/paper_033_execution_receipt.json}{execution
receipt}, \href{code/cross_domain_symbol_game_receipt.json}{cross-domain
receipt}, and earlier release validation records are historical
artifacts until their source and executable bindings are reconciled with
a new run. A SHA-256 digest can bind bytes. It does not establish their
truth, authenticate an unsigned producer, or supply a missing
experiment. A restricted sorted-JSON encoding also needs its own stated
contract; it should not be called full
\href{https://www.rfc-editor.org/rfc/rfc8785}{RFC 8785} conformance
without checking that \nolinkurl{specification.} Run tests and lint
separately from the default deposit verifier: the latter checks byte
custody and structure rather than executing the entire suite or proving
the answers.

\subsection{7. Map a new problem
carefully}\label{map-a-new-problem-carefully}

Before adapting the example, write down what a vertex means, what an
edge value measures, why the edge-difference relation is appropriate,
and which changes are allowed. Supply units and direction conventions.
Decide whether changing a claim is an acceptable correction or would
replace an observation you must preserve.

For example, measured offsets might need an uncertainty model rather
than unrestricted replacement. A dependency relation might be Boolean
rather than an additive quantity. Encoding either as this graph is a
modeling step that requires \nolinkurl{justification.} Renaming vertices
to ``currency,'' ``permission,'' or ``quantum state'' does not supply
that \nolinkurl{justification.}

The package's \href{code/multi_domain_symbol_game_bouts.py}{four labeled
graph scenarios} illustrate repeated use of the finite interface. Their
labels do not establish market equilibrium, authorization security, or a
quantum no-go theorem. The governed scripts make no LLM calls, so
hypothetical unconstrained-chat costs in older records are not a
measured comparison.

For model integration, keep the proposer behind the same boundary: parse
the proposed fields, check the declared carrier and operations,
recompute the relevant property, and save the actual response. Consult
\href{code/llm_streaming_harness.py}{the streaming harness} as an
implementation example; reviewing its adapter does not establish a live
provider run. The harness accepts integers and exact rational strings as
claim values and refuses a JSON float or Boolean before any scoring,
since \nolinkurl{Fraction(0.1)} is the binary expansion of the float
rather than one tenth
(\nolinkurl{code/tests/test_harness_exact_inputs.py}).

\subsection{8. Connections and next
reading}\label{connections-and-next-reading}

The familiar idea of an untrusted producer supplying evidence to a
checker appears in Necula's
\href{https://doi.org/10.1145/263699.263712}{proof-carrying code}. That
work concerns proofs of a declared code-safety policy. It supplies
context for separating proposal production and evidence checking; this
tutorial does not claim its graph receipt is such a proof.

The minimization and lower-bound reasoning connect to linear-programming
duality; see Boyd and Vandenberghe,
\href{https://doi.org/10.1017/CBO9780511804441}{Convex Optimization},
Chapter 5. The series'
\href{https://github.com/dbccpage/the_symbol_game/blob/59dde95a7860063c25d75db37f40af556e9f28ee/omega_engine_atlas/coda_atlas/F004_papers/015_pfp/001_pfp_discrete_cochain_retraction_and_toy_involutions/paper.md}{graph
mathematics paper} develops the finite theory;
\href{../004/paper.pdf}{Paper 004} explains native verification and
evidence custody. The
\href{https://github.com/dbccpage/the_symbol_game/blob/59dde95a7860063c25d75db37f40af556e9f28ee/omega_engine_atlas/coda_atlas/F004_papers/015_pfp/007_pfp_executable_lessons_and_evidence_boundaries/paper.md}{curriculum-design
guide} turns examples into lessons, and the
\href{../009/paper.pdf}{review protocol} explains how to audit a claimed
result.

This guide supports a finite, inspectable proposal--check workflow. Its
practical success criterion is that the reader can state the original
instance, explain the failed proposal, distinguish residual
concentration from claim replacement, and identify exactly what a fresh
receipt records. General semantic correctness and model quality need
their own evidence.

\subsection{9. Limitations}\label{limitations}

The guide works one three-edge instance by hand and describes one
scripted session. The current publication review checks the session and
the new input-boundary cases on the exact Python route. Native execution
was unavailable in that review; native-dependent tests are explicitly
skipped, and a mocked native error test checks propagation and rollback
without claiming to run Rust. The fresh
\href{docs/publication_review_2026_10_02.json}{review record} binds
commands, outputs, tools, source comparisons, and tested files.

The earlier \href{paper_change_log.md}{author-review history} reports a
Windows native run and 75 passing package tests; its executable digest
remains in the historical records and claim ledger. That report concerns
earlier bytes and is not execution evidence for this patched package.
The executable is not deposited. A current native acceptance claim needs
a fresh build, explicit pins, and a run of the patched suite. The legacy
scenario and streaming scripts remain implementation examples, not
measured comparisons. The examination's UNKNOWN budget outcome is an
instructional policy for a separately declared bounded search; the
four-move governor does not implement that terminal status or budget
controller.

This paper assigns no authorization and changes no
\nolinkurl{jurisdiction.} It carries no operational authority or
institutional standing.

\subsection{Verification for this
revision}\label{verification-for-this-revision}

The package verifier (\nolinkurl{verify_deposit.py}) checks the manifest
and two-way file census. The \href{tools/rebuild_release.py}{release
tool}, run with \passthrough{\lstinline!--check!}, checks the archive
against the manifests. The
\href{docs/publication_review_2026_10_02.md}{publication review}
describes the corrections and remaining replay boundary. Its separate
JSON execution record supplies the current test and lint evidence.
Earlier receipts retain their dates and original scope. Version 0.1.0
remains the initial preprint candidate; no external deposit or DOI is
claimed.

\subsection{Practitioner extension: run a bounded harness
experiment}\label{practitioner-extension-run-a-bounded-harness-experiment}

Freeze the source, checker, model configuration,
\nolinkurl{public/private} task partitions and original budget before
proposing rewrites. Use the research command's
\passthrough{\lstinline!preflight!} operation to prepare the protocol
and inspect the report. The offline form spends nothing. Live preflight
additionally requires the verified Ubuntu host, a pinned local container
image, current pricing and provider \nolinkurl{compatibility.} The
reproduction guide in Paper 029 gives the complete commands.

Use \passthrough{\lstinline!run!} for the experimental sequence and
\passthrough{\lstinline!resume!} for an interrupted sequence with the
same retained directory. An unreconciled provider reservation is a
reason to stop dispatch, not to create a fresh budget directory. Inspect
the raw response, usage observation and attempt identity before
reconciling it. Never treat a failed model answer as a free call merely
because it supplied no usable program.

Read public completion and private completion as different columns. Both
selection and confirmation are optimization data. Only the untouched
transfer panel supports the predeclared transfer comparison. A failed
proposal and an unchanged reference incumbent are valid outcomes to
retain. Use \passthrough{\lstinline!replay!} and
\passthrough{\lstinline!report!} to reconstruct results offline; replay
does not buy another independent model trial. This revision has no live
campaign results and supplies no basis for a performance
\nolinkurl{recommendation.}

The source and page references for this extension are collected in
\href{../029/paper.pdf}{Paper 029} and its
\href{../029/docs/source_to_claim_map.md}{source-to-claim map}. The
author-review rendering does not replace a historical deposit or
archive.

\subsection{References}\label{references}

\begin{itemize}
\tightlist
\item
  Necula, G. C. (1997).
  \href{https://doi.org/10.1145/263699.263712}{Proof-carrying code}. In
  \emph{Proceedings of the 24th ACM SIGPLAN-SIGACT Symposium on
  Principles of Programming Languages (POPL '97)}, pp.~106--119. ACM.
  DOI: \nolinkurl{10.1145/263699.263712.}
\item
  Boyd, S., and Vandenberghe, L. (2004).
  \href{https://doi.org/10.1017/CBO9780511804441}{Convex Optimization}.
  Cambridge University Press, Chapter 5. DOI:
  \nolinkurl{10.1017/CBO9780511804441.}
\item
  Rundgren, A., Jordan, B., and Erdtman, S. (2020).
  \href{https://www.rfc-editor.org/rfc/rfc8785}{JSON Canonicalization
  Scheme (JCS)}. RFC 8785 (Informational), RFC Editor, June 2020. DOI:
  \nolinkurl{10.17487/RFC8785.}
\end{itemize}

\end{document}
